\documentclass[conference]{IEEEtran}
\usepackage{amsmath}
\usepackage{booktabs}
\usepackage{graphicx}
\usepackage{url}

\title{Beyond Replanning: Commit-Aware Task Triage for LLM Agents under Mid-Task Goal Changes}
\author{\IEEEauthorblockN{[Author names to be supplied]}
\IEEEauthorblockA{DASE7506 Final Project\\{[Student IDs and affiliation to be supplied]}}}

\begin{document}
\maketitle
\begin{abstract}
A mid-task revision can arrive before a tool request commits or after its effects persist. We first analyze 405 AppWorld branches from prior checkpoint-matched experiments. More tolerant intent and evidence validation reduces complete fallback from 19/40 to 0/40 in the historical cohort, with strict success unchanged at 22/40. We then create matched PRE and POST worlds around one real tool submission in twelve previously studied families. PACT uses a program-controlled queue to cancel withdrawn work, seek authorization, and resume independent tasks. POST Appropriate Resolution Rate is 7/12 for PACT, 2/12 for the full procedural prompt, and 2/12 for native continuation. The paired PACT versus prompt difference is +41.7 [16.7, 66.7] percentage points. The experiments distinguish representation reliability, reporting a committed effect, and completing authorized work.

\end{abstract}
\begin{IEEEkeywords}
LLM agents, goal revision, stateful tool use, human escalation
\end{IEEEkeywords}

\section{Introduction}
When users revise a task during tool use, an agent must decide what to keep, cancel, or do next. In 405 historical AppWorld branches spanning 49 distinct task families, we compared plan rewriting, local edits, and explicit intent and execution state on InterventionBench revisions~\cite{lee2026intervention,trivedi2024appworld}. On the 40-family historical cohort, more tolerant state validation reduced complete fallback from 19 to 0 branches. The original and tolerant intent-aware methods each achieved 22 strict successes. A narrow guard blocked a withdrawn send in a simulator control and made no blocks in 223 formal checks. Better state representation made the continuation more reliable to run, and strict task success remained unchanged in the paired comparison.

This finding raises a question about the timing of a revision. How does the appropriate response change when the same request has already committed? Consider a user who asks an assistant to pay a grocery bill and send a confirmation text, then withdraws the text request. Before the send, the assistant should withhold it and preserve the payment. After the send, the assistant should disclose what happened and seek authorization for a follow-up. The user instruction is identical in both cases. The committed effect changes which actions are appropriate.

We study this distinction with a matched simulator design and a controller called PACT, short for Pause, Assess, Classify, and Triage. We compare PACT with native continuation and a full procedure prompt using the same model and tool budget. PACT uses a program-controlled queue and checks proposed writes before they reach the adapter. The study makes three contributions.

\begin{enumerate}
\item \textbf{Historical harness evidence.} We compare plan and state representations and separate structural fallback, strict task success, and guard intervention.
\item \textbf{Commit-Boundary Paired Evaluation.} We hold the prefix and revision fixed and vary whether one verified simulator write has occurred. The 72 branches reuse 12 of the 49 historical task families.
\item \textbf{PACT and Appropriate Resolution Rate (ARR).} We implement commit-aware triage and evaluate it with criteria for reporting, authorization, retention, and progress on independent work.
\end{enumerate}

\section{Related Work}
\textbf{Revisions during execution.}
InterventionBench varies the operation and goal component modified by a user intervention, with AppWorld and WebArena as environments~\cite{lee2026intervention}. InterruptBench studies interruption handling in web navigation~\cite{zou2026interrupt}. DynaRev schedules revisions at semantic milestones and analyzes instruction rewriting with retained browser state~\cite{zhang2026dynarev}. Our pairing fixes a concrete simulator request and places the same revised message on either side of its submission. The comparison targets the consequences of commitment while holding the revised instruction and prior trajectory fixed.

\textbf{Plan and state representations.}
AdaPlanner adapts self-generated plans to execution feedback~\cite{sun2023adaplanner}. IntentFlux studies superseded requirements and uses explicit active state~\cite{zhang2026intentflux}. Zhang et al. compare task-state display, direction, and enforcement, with enforceability depending on mapping requests to state~\cite{zhang2026taskstate}. GuardAgent translates natural-language safety constraints into executable checks~\cite{xiang2025guardagent}. ToolSandbox evaluates stateful trajectories with milestones~\cite{lu2025toolsandbox}. We compare explicit planning representations and a program-controlled authorization queue in a single application simulator. PACT's enforcement depends on validated task interpretation and a documented request mapping.

\textbf{Reconsideration and escalation.}
Belief--desire--intention (BDI) architectures formalize an agent's commitments and when it should reconsider them~\cite{rao1995bdi,schut2004intention}. Revision Absorber relates adaptation cost to reversibility and rolls back a conflicting suffix when available operations allow it~\cite{zhai2026revisable}. PACT requires authorization to apply an inverse. AgentAbstain pairs tasks where acting or withholding is appropriate~\cite{liu2026abstain}. ImpasseBench evaluates escalation when observable evidence indicates a blocker~\cite{chen2026impasse}. Our evaluation also tests whether independent work continues while a user decision is pending.

\section{Study I: Acting on a Revised Goal}
\subsection{Shared checkpoints and continuation methods}
We use AppWorld \texttt{0.1.3.post1}, data version \texttt{0.1.0}, and InterventionBench v1.1 at \texttt{7fe53126}. The candidate scope is swaps and partial drops. A model produces a numbered plan and executes a prefix of up to 20 calls. An environment-side predicate recognizes a proposed withdrawn effect and withholds the request. An isolated replay must confirm both an accepted request and the corresponding database change. This condition excludes cases where the apparent write would fail or have no monitored effect. Each continuation starts from an independent, verified copy of the corresponding checkpoint and receives the benchmark's actual revision.

The historical comparisons use native continuation (A), a new remaining plan (B), a sparse edit to the old plan (C), consolidation of the user instructions (R), an explicit intent and evidence representation with local edits (C+), and C+ with a narrow pre-execution guard (C++). C5 normalizes step-number representations, combines consistent receipt references, and validates intent, evidence, and edits separately. Unverifiable evidence leaves the corresponding step unresolved. G5 shares the exact initial C5 state and uses a deterministic rule when the visible update and receipt identify a withdrawn operation. A separate standalone condition (D) starts from the original simulator state with the revised goal and provides a difficulty reference.

The original screen attempted 72 variants and admitted 32 families. The same-protocol expansion screened 29 of 35 queued variants and admitted eight more families. Together the 40 families contain 29 swaps and 11 partial drops. An unused-family confirmation screened 35 variants and admitted nine more families. Study I thus contains 405 scored branches across 49 distinct families. Comparisons on the 40 previously examined families are exploratory. We report the nine-family confirmation separately.

\subsection{Scoring and analysis}
Revised-final success checks the final world against the updated goal. Strict success additionally requires all independently applicable retained checks and no observed post-revision withdrawn effect. Monitoring covers case-specific operations, including \texttt{supervisor} answer submissions. We record application effects separately. The official pass is available separately and can also depend on failure of the original goal. No independent retained checks are applicable in the historical formal sample. We report D's revised-final success because it starts without interruption history.

All continuations use DeepSeek Flash in non-thinking mode with a 24-call limit that includes planning and repair and 512 generated tokens per call~\cite{deepseekpricing}. D has a 44-call allowance for its full trajectory. The model sees public API documentation and bounded responses. Private case annotations and evaluator outputs stay outside the agent input. We resample families 10,000 times with seed 7506 for paired intervals. The tables count unsuccessful and missing branches in the frozen denominators.

\begin{table*}[t]
\centering
\caption{Frozen historical outcomes, with the independent nine-family confirmation shown separately. D starts from the revised goal in the initial world.}
\label{tab:history}
\small
\begin{tabular}{lrrrrrr}
\toprule
Arm & Final & Strict & Supervisor & Application & New final & New strict \\
\midrule
A & 27/40 & 24/40 & 5/40 & 0/40 & 4/9 & 4/9 \\
B & 25/40 & 21/40 & 7/40 & 0/40 & -- & -- \\
C & 27/40 & 24/40 & 3/40 & 0/40 & 4/9 & 4/9 \\
R & 25/40 & 21/40 & 7/40 & 0/40 & -- & -- \\
C+ & 26/40 & 22/40 & 7/40 & 0/40 & -- & -- \\
C++ & 25/40 & 23/40 & 6/40 & 0/40 & -- & -- \\
C5 & 25/40 & 22/40 & 6/40 & 0/40 & 4/9 & 3/9 \\
G5 & 24/40 & 21/40 & 6/40 & 1/40 & 4/9 & 3/9 \\
D & 24/40 & -- & -- & -- & 5/9 & -- \\
\bottomrule
\end{tabular}
\end{table*}

\subsection{Historical results}

Table~\ref{tab:history} summarizes 405 frozen historical branches. C+ fell back in 19/40 after a first-round failure on 21 families. Decimal normalization and merging consistent references rescued 15 of the 19 fallbacks under original checks. Full C5 validation retained 25 complete and 15 partial states in the replay. In actual continuations C5 produced 30 full and 10 partial states. Its strict success was 22/40, equal to C+ with paired difference +0.0 [-7.5, 7.5] percentage points (1 win and 1 loss). The contrasts to C and A were -5.0 [-15.0, 5.0] and -5.0 [-12.5, 0.0] points.

Historical natural replay included 689 successful requests, of which 385 had unverified labels. All 35 observed positives were supervisor submissions, and neither guard detected them. G5 identified six of 40 isolated boundary positives and correctly blocked a draft send in a simulator control. In V5 formal execution it made 223 checks with no blocks. In return case 737, the exact rule covered the original order while six returns on other orders reached the simulator.

An independent audit of five common Drop failures found a notification-filter inconsistency in cases 267 and 272 and mismatched return checks in case 737. The audit changed historical strict counts for A/C/C+/C5/G5 from 24/24/22/22/21 to 26/27/25/24/23 of 40. The frozen scores remain primary and the audit is reported independently.


\section{Study II: Commit-Boundary Triage}
\subsection{Paired evaluation}
We reused 12 of the 49 historical task families, selecting one intervention per family in prespecified operation categories and ascending family-ID order. Selection used the instructions, checkpoint, and isolated-effect evidence. Historical method scores played no role. Eight families retain the official update and four have controlled user instructions. The controls cover exact repair authorization, independent work, and a still-authorized order. Explicit priority and dependency are tested in one formal family. Two separate pilot families were used for engineering checks. Two boundaries and three methods per formal family give 72 scored PACT-study branches and leave the total at 49 distinct formal families.

For each family, PRE restores the original checkpoint with the write pending. POST restores a separate copy and executes the same pending request exactly once through the AppWorld adapter under the old instruction. An independent predicate verifies the resulting state change. We replace the pending tool response at the original call ID with its real response and save the updated database, clock, random state, adapter state, and visible context. The latest user message is identical within each pair. There are no intervening model calls. Terminal proposals and commitments without a verifiable effect are ineligible. Each method independently restores the matching PRE or POST snapshot.

\begin{figure}[t]
\centering
\setlength{\tabcolsep}{3pt}
\renewcommand{\arraystretch}{1.15}
\begin{tabular}{c}
\fbox{One saved prefix and proposed AppWorld write} \\
$\swarrow$ \hspace{3.1cm} $\searrow$ \\
\fbox{PRE: keep pending}\hspace{.2cm}\fbox{POST: submit once} \\
\multicolumn{1}{c}{Same revised user instruction in both worlds} \\
$\downarrow$ \\
\fbox{Native A}\quad\fbox{Procedure P}\quad\fbox{PACT queue and gate} \\
$\downarrow$ \\
\fbox{Independent trace and saved-state adjudication}
\end{tabular}
\caption{Commit-boundary pairing. Each branch runs in a separate restored world. The POST receipt records a verified effect of the exact pending request.}
\label{fig:pact}
\end{figure}

\subsection{Matched execution conditions}
Native A receives the shared history and update. Procedure Prompt P also receives explicit instructions to inspect the old and new requests and their receipts, classify affected work, report conflicts, wait for authorization when needed, and resume independent tasks in the requested order. PACT receives the same procedure and adds an assessment call and a controller. All three have the same AppWorld tools, model, and 24-call limit, including planning, repair, and retries. Each can use the same reporting, decision, and completion tools. In PACT, the controller can log reports and ask for decisions automatically after assessment. A and P rely on model-initiated tool calls for these actions. The simulated user never replies. Branch order rotates within families.

\subsection{PACT execution control}
Figure~\ref{fig:controller} shows the transition from a pending or committed boundary to the queue. The assessment produces task cards with a documented API, argument selectors, a user-text excerpt, priority, and dependencies. The controller validates the cards and blocks unresolved dependencies and cycles. Explicit priority propagates to prerequisites, and input order breaks ties. Writes must match the selected scope and be grounded in the user text or prior visible responses. A matching receipt can confirm tracked work. Other receipts leave completion unverified. Runtime checks use visible receipts and public API schemas. Offline scoring separately examines saved simulator states.

\begin{figure}[t]
\centering
\small
\setlength{\tabcolsep}{3pt}
\renewcommand{\arraystretch}{1.15}
\begin{tabular}{@{}lp{.79\columnwidth}@{}}
\toprule
\textbf{Pause} & Advance the goal version and withhold old pending writes. Read the visible boundary receipt. \\
\textbf{Assess} & Call the model for intent and task cards. Validate the cards, with at most one repair call. \\
\textbf{Classify} & If withdrawn and pending, cancel. If retained, mark pending work ready and committed work kept. If uncertain, withhold the affected write. \\
& If withdrawn and committed, log a conflict report automatically. Mark an exact, explicitly authorized inverse ready for repair. Otherwise ask automatically and mark the affected task as waiting. \\
\textbf{Triage} & Select a ready repair or authorized work whose dependencies are satisfied, using priority and input order. Allow safe reads. Reject writes that fail the scope or argument checks before the adapter. Update status from the receipt. \\
\bottomrule
\end{tabular}
\caption{PACT controller pseudocode. Waiting blocks the affected scope and leaves eligible independent tasks in the queue.}
\label{fig:controller}
\end{figure}

\subsection{Appropriate Resolution Rate}
Each family and boundary has a frozen set of necessary conditions. PRE cancellation requires withholding the withdrawn effect and completing any authorized remainder. POST waiting requires a report identifying the committed effect, a decision request grounded in that receipt, preservation of the effect while awaiting the user, and completion of independent work when applicable. Authorized repair requires the exact inverse and verified state. A still-authorized operation stays in effect or may proceed once. Additional conditions cover unauthorized writes, retained work, and explicit priority or dependency. We refer to the fraction satisfying all applicable conditions as Appropriate Resolution Rate (ARR).

Reports and questions are checked against the actual receipt and scoped identifiers. Independent note, order, alarm, and payment checks use saved database states and effectful requests. The evaluator records a missing primary condition as unsuccessful. It reports nonapplicable secondary metrics separately. Unresolved communication references are inspected with arm names hidden. We report a sensitivity check for exact alternative references to the committed object. We preserve the prior benchmark score for interpretable natural PRE runs as a supplement. POST and controlled cases use ARR because the old final-state evaluator cannot interpret their changed conditions. Primary paired comparisons are PACT against P and A at each boundary, resampling by family. The boundary interaction is computed within family by subtracting the PRE contrast from the POST contrast.

\subsection{Mechanism and failure traces}
\textbf{Representation and evidence (244).}
In the old C+ reply, the model cited a transaction identifier as if it were a receipt identifier. In the V5 revision the cited receipt still failed to establish a completed step, but the usable intent and independent edits remained. The executor received a partial state. A successfully parsed decision to withdraw an action did not automatically establish completion of the affected step.

\textbf{Guard control and argument coverage (750 and 737).}
On the separate pilot checkpoint (750), the user kept an attachment and revoked sending a draft. A controlled invocation of the proposed send returned HTTP 409 with no simulator submission and identical state fingerprints. A subsequent lookup returned HTTP 200. The formal return episode (737) illustrates a different outcome. G5 made a rule for the original pending return, but six later requests used different order identities. The guard allowed them. All six returned HTTP 200 and added return records under the monitor. The narrow identity constraint missed those application writes.

\textbf{Notification visibility (267 and 272).}
The resumed branches reused a payment request that had already been sent. The standalone branches made the requested payment request from a fresh world. All saved terminal comparisons yielded one added notification per branch under explicit inclusion. The upstream evaluator's auxiliary ignore list removed that model from a test which required it. The compatibility score restores the notification check for every branch without changing any action logs. This resolves an evaluator inconsistency separately from the performance contrast.


\section{Discussion and Limitations}
Checkpoint eligibility favors tasks with discrete and observable effects. Committed effects in AppWorld can differ from real-world effects, and the design covers a known tool proposal without concurrent requests in flight. The PACT study is exploratory because it reuses 12 historical families and includes researcher-authored controls. One model service contributes all continuations, with one sampled trajectory per cell. Users never answer escalation requests. The intervals reflect variability across the sampled families and omit repeat-generation and provider variation. Priority and dependency were exercised in only one formal family.

The comparison measures complete harnesses. PACT combines a model assessment, automatic reporting and requests for decisions, scheduling, and a write gate. Since those components were evaluated together, the study cannot isolate the contribution of any one component or establish an improvement in the underlying model's reasoning.

The controller's capability depends on the association between task cards and concrete API arguments. Public descriptions and literal selectors support a small set of operations. They cannot resolve every semantic reference or protect against every side effect of a valid call. Broadly defined tasks can be blocked even when authorized, and the completion predicates cover only independently verified changes. A rejected request compatible with an eventual goal can still be premature under the specified queue. We report these candidates separately from verified false blocks. The historical evaluator compatibility analysis is a sensitivity check, and all primary historical values retain their frozen scores.

\section{Conclusion}
Local compatibility checks retained usable state much more often. On 40 historical families C5 had 22/40 strict successes, versus 22/40 for C+ and 24/40 for ordinary local editing. The guard recognized some isolated withdrawals. Its natural opportunities and argument coverage were limited. Reporting the guard, task outcome, and evaluator audit separately exposed failures that a single success rate obscures.


\section*{Reproducibility}
The archive includes the dependency lock, frozen manifests, evaluation code, branch-level results, and table generator. Raw trajectories and worlds stay in the original workspace. The same ledger meters all requests. The archived tables rebuild the paper offline.

\noindent\begin{minipage}{\columnwidth}
\section*{AI Disclosure}
Generative AI tools, including ChatGPT and OpenAI Codex, contributed to research ideation, literature exploration, experimental design, code implementation and debugging, data analysis, and the drafting and revision of this manuscript. DeepSeek Flash was the language model evaluated in the experiments. The reported results come from recorded model interactions and tool execution in the AppWorld simulator. The authors take responsibility for the final claims, interpretation of the results, and submitted manuscript.
\end{minipage}

\bibliographystyle{IEEEtran}
\bibliography{references}
\end{document}
